\documentclass[11pt,a4paper]{article}
\usepackage[T1]{fontenc}
\usepackage{lmodern,amsmath,amssymb,amsthm,mathtools,booktabs,microtype}
\usepackage[margin=27mm]{geometry}
\usepackage[colorlinks=true,linkcolor=blue,urlcolor=blue,citecolor=blue]{hyperref}
\newtheorem{theorem}{Theorem}[section]
\newtheorem{lemma}[theorem]{Lemma}
\newtheorem{corollary}[theorem]{Corollary}
\newtheorem{proposition}[theorem]{Proposition}
\newcommand{\Q}{\mathbb Q}\newcommand{\Z}{\mathbb Z}
\DeclareMathOperator{\Res}{Res}
\DeclareMathOperator{\Norm}{Norm}
\DeclareMathOperator{\Disc}{Disc}
\setlength{\emergencystretch}{2em}
\title{Separating the Nonzero Part of a Finite Algebra}
\author{}\date{9 November 2026}
\begin{document}\maketitle
\begin{abstract}
We give an explicit polynomial, using one additional rational coordinate,
for a guarded rational point of a finite algebra obtained by adjoining
quadratic roots to an algebra with separable fibers. A characteristic
polynomial identifies the part on which the guard is nonzero.
Norms of discriminants then give a rational separating element by a
root-separation estimate. The construction covers degree drops,
nonreal points, and nonreduced fibers where the guard vanishes.
Applied to the specified arithmetic sources, it gives a family in
twenty-seven total variables of degree less than $2^{109}$ and a
fully specified polynomial in twenty-nine total variables of degree
less than $2^{111}$ and coefficient height at most $2^{2^{121}}$.
The algebraic elimination theorem is unconditional. The assertions
about defining the integers retain the arithmetic premises and the
effective or ineffective thresholds of the respective sources.
\end{abstract}

\section{The algebra and its guarded rational points}

Let $R_0=\Z[x_1,\ldots,x_a]$. Consider the finite free algebras
\[
 B=R_0[z_1,\ldots,z_p]/(f_1(z_1),\ldots,f_p(z_p)),\qquad
 A=B[q_1,\ldots,q_s]/(q_j^2-c_j)_{j=1}^s ,
\]
where $p,s\ge1$, each $f_i\in R_0[Z]$ is monic of degree
$d_i\ge2$, and $c_j\in B$.
Assume that every $f_i$ is separable after every rational
specialization of the base coordinates. Put $m=p+s$ and write
$\theta_1,\ldots,\theta_m$ for the generators in this order,
with native degrees $d_{p+j}=2$. Let
\[
 b=\prod_{i=1}^p d_i,\qquad n=b2^s.
\]
The native monomials, with each exponent smaller than its defining
degree, are bases of $B$ and $A$. For $u$ in either algebra,
$M_u$ denotes its multiplication matrix in the relevant basis.
Its determinant is the algebra norm.

Fix $g\in B$ and the guard
\[
                  G=gq_1\cdots q_s,\qquad
                  T=g^2c_1\cdots c_s\in B.
\]
Thus $G^2=T$ in $A$. The predicate to be encoded at a rational
base point is
\begin{equation}\label{eq:predicate}
 \text{there exists a homomorphism }A\longrightarrow\Q
 \text{ over that base point with }G\ne0.
\end{equation}
All finite algebra generators will be eliminated; exactly one
new rational coordinate $U$ is used. There are no further residual
equations in the predicate.

At a rational base point, $B$ is an \'etale $\Q$-algebra.
On each geometric factor where $T\ne0$, all $c_j$ and $g$ are
nonzero. The corresponding part of $A$ is therefore \'etale.
The other factors of $A$ may be nonreduced; they are retained
throughout the construction.

\section{An explicit polynomial}

Define
\[
 \det_B(I-vM_T)=\sum_{j=0}^b a_jv^j,\qquad a_0=1.
\]
The native discriminants are elements of $B$:
\[
 \delta_i=\Disc(f_i)\quad (i\le p),\qquad
 \delta_{p+j}=4c_j.
\]
Use the usual discriminant convention
$\Disc(f)=(-1)^{d(d-1)/2}\Res(f,f')$ for a monic polynomial.
Set
\begin{equation}\label{eq:root-bound}
 R=n\left(1+\sum_{i=1}^m\sum_{u,v=1}^n
                         (M_{\theta_i})_{uv}^{\,2}\right),
 \qquad S=b\sum_{i=1}^m d_i(d_i-1).
\end{equation}
These are an integer polynomial and a positive integer,
respectively.

For each $r=1,\ldots,b$, perform the following finite polynomial
operations:
\begin{align*}
 c&=a_r,& h_r(Z)&=\sum_{j=0}^r a_jZ^{r-j},&
 E&=c-h_r(T)\in B,\\
 D_i&=\Norm_{B/R_0}(E\delta_i+c-E),&
 D&=\prod_{i=1}^mD_i^2,\\
 K&=(2R)^S(1+c^2)^{mb}\prod_{i=1}^m(1+D_i^2),&
 C&=D+4RK.
\end{align*}
The symbols $c,E,D_i,D,K,C$ depend on $r$.
Define in $A$
\begin{equation}\label{eq:beta}
 \beta_r=cDE\left(C^m+
                    \sum_{i=1}^mD^iC^{m-i}\theta_i\right),
\end{equation}
and finally define integer polynomials
\begin{align}
 H_r(\mathbf x,U)&=\det_A(I-UM_{\beta_r}),\label{eq:H}\\
 Q_r(\mathbf x,U)&=H_r(\mathbf x,U)^2+
                         \sum_{j=r+1}^b a_j(\mathbf x)^2,\label{eq:Q}\\
 P(\mathbf x,U)&=\prod_{r=1}^b Q_r(\mathbf x,U).\label{eq:P}
\end{align}
Every matrix and finite index range is specified; no division
occurs in these formulas. In particular they give an integer
polynomial even where $cD=0$.

\begin{theorem}\label{thm:main}
For every rational base point $\mathbf x$, predicate
\eqref{eq:predicate} holds if and only if there is a rational
$U$ with $P(\mathbf x,U)=0$.
\end{theorem}

\section{Proof of the projection theorem}

Fix a rational base point for this section. A geometric character
of an algebra means a homomorphism from that algebra to
$\overline{\Q}$ over $\Q$. Call a character of $B$ active when
$T$ is nonzero there. A character of $A$ is active when its
restriction to $B$ is active; this is equivalent to $G\ne0$.

\begin{lemma}\label{lem:idempotent}
The degree of $\det_B(I-vM_T)$ is the number $r_0$ of active
geometric characters of $B$, counted individually.
If $r=r_0>0$, then $c\ne0$ and $e=E/c$ is the idempotent
equal to one on the active factors and zero on all other factors.
\end{lemma}
\begin{proof}
Since $B$ is \'etale, $M_T$ is diagonalizable over
$\overline{\Q}$. Its determinant is the product of
$1-vT(\chi)$ over all $b$ characters $\chi$.
For $r=r_0$, the polynomial $h_r(Z)$ is the product of
$Z-T(\chi)$ over the active characters, including repetitions
of their nonzero values. It vanishes at every such value and
has constant term $c$. Thus $E=c-h_r(T)$ has value $c$ on
every active factor and zero on every inactive factor.
\end{proof}

For the correct case $r=r_0>0$, the element
$e\delta_i+1-e$ is a unit of $B$. Indeed the point discriminants
never vanish and the pole discriminants do not vanish on active
factors. Hence
\begin{equation}\label{eq:Di}
 D_i=c^b\prod_{\chi\ {\rm active}}\delta_i(\chi)\ne0.
\end{equation}
In particular $D>0$, $K>0$, and
\[
                         \varepsilon=D/K\in(0,1].
\]

\begin{lemma}\label{lem:separate}
In the correct case, if two active characters of $A$ first
differ at $\theta_i$, their values there differ in complex
absolute value by at least $\varepsilon$. The element
\[
 \alpha=L^m+\sum_{i=1}^mL^{m-i}\theta_i,
                       \qquad L=C/D=1+4R/\varepsilon,
\]
has nonzero, pairwise distinct values at the active characters.
\end{lemma}
\begin{proof}
Every value of every native generator is an eigenvalue of its
multiplication matrix and has absolute value at most $R$.
For example, the spectral radius is at most $n$ times the
largest absolute entry, and $1+t^2\ge |t|$ bounds each entry
by the parenthesis in \eqref{eq:root-bound}.

Let $\Delta$ be the difference at the first differing generator.
For a point generator, its defining polynomial has scalar
coefficients. For a pole generator, the earlier point
coordinates coincide, so both values are roots of the same
quadratic. In either case $\Delta^2$ is one of the squared
differences occurring in the native discriminant at an active
character of $B$. Put $k_i=bd_i(d_i-1)$ and
\[
 z=\left|\prod_{\chi\ {\rm active}}\delta_i(\chi)\right|
                         =|D_i|/|c|^b.
\]
All other root differences have absolute value at most $2R$.
There are at most $k_i$ differences when squared factors are
counted twice. Consequently
\[
                  |\Delta|\ge
             \frac{\sqrt z}{(2R)^{(k_i-2)/2}}.
\]
On the other hand,
\[
 \frac{D_i^2}{(1+D_i^2)(1+c^2)^b(2R)^{k_i}}
 \le\frac{\min(1,z^2)}{(2R)^{k_i}}
 \le\frac{\sqrt z}{(2R)^{(k_i-2)/2}}.
\]
Each left-hand factor is at most one, and their product over
$i$ is $\varepsilon$. This proves the required lower bound.
The argument allows complex roots and repeated values of $T$.

If the first differing index is $i$, the leading term in the
difference of the two $\alpha$-values has absolute value at
least $\varepsilon L^{m-i}$. The remaining terms have sum
of absolute values less than
\[
       \frac{2R}{L-1}L^{m-i}
                       =\frac{\varepsilon}{2}L^{m-i}.
\]
The values are therefore distinct. Similarly, the sum after
the constant $L^m$ has absolute value less than
$RL^m/(L-1)\le L^m/4$, so no value is zero.
\end{proof}

\begin{proof}[Proof of Theorem~\ref{thm:main}]
If $c=0$ or $D=0$ in a case, then $\beta_r=0$ and $H_r=1$.
If $Q_r=0$ at rational coordinates, its sum-of-squares form
therefore gives $cD\ne0$ and $a_j=0$ for every $j>r$.
Thus this is precisely the correct positive degree case.

By Lemma~\ref{lem:idempotent}, multiplication by $\beta_r$
vanishes identically on the inactive part of $A$; this
includes its nilpotents. On the active part,
\[
                       \beta_r=c^2D^{m+1}\alpha.
\]
That part is \'etale, and Lemma~\ref{lem:separate} shows that
these values are distinct and nonzero. Hence $H_r$ is the
product of distinct factors $1-U\beta_r(\chi)$ over the active
characters. A rational root gives a one-dimensional rational
kernel of $I-UM_{\beta_r}$. Every native multiplication matrix
preserves this kernel and acts on it by a rational scalar.
Those scalars satisfy all the algebra relations. The resulting
rational character lies on the active part, since the
$\beta_r$-eigenvalue is nonzero. This proves soundness.

Conversely, a guarded rational character implies $r_0>0$.
In the correct case $cD\ne0$, and its $\beta_r$-value is a
nonzero rational number. Its reciprocal annihilates $H_r$;
all tail coefficients vanish, so it annihilates $Q_r$ and
$P$. If $r_0=0$, every $a_r$ is zero, every $H_r$ is one,
and $P$ has no zero. This also covers the empty active part.
\end{proof}

\section{Degree and coefficient estimates}

Give the base variables weight one and the native generators
weights $w_i\ge1$, such that monic reduction never increases
weight. Use a weighted coefficient column norm on multiplication
matrices: each scalar polynomial has its coefficient $\ell^1$
norm, and the basis weights for that norm are
$2^{\sum a_i\gamma_i}$ for the native monomial
$\prod_i\theta_i^{a_i}$. Suppose
\[
 \|M_{\theta_i}\|\le2^{\gamma_i},\qquad
 \operatorname{wt}(G)\le w_G,\qquad \|M_G\|\le2^\eta.
\]
All logarithmic bounds below are nonnegative integers.
The companion-column recurrence in \cite{Weighted}
provides these data directly from the integer coefficient circuit.
The norm and weight restricted to $B$ use its first $p$ factors.
Because $G=gq_1\cdots q_s$, substitution of $q_j^2=c_j$
gives the bounds $t=2w_G$ and $\mu=2\eta$ for $T$ in $B$.
For the norm assertion, $B$ is an invariant subspace for
$M_{G^2}$ and its basis weights are the restricted weights.

For clarity write $\ell(u)=\lceil\log_2u\rceil$ for positive
integers $u$. Set
\begin{align*}
 w_*&=\sum_i(d_i-1)w_i+\max_i w_i,&
 \gamma_*&=\sum_i(d_i-1)\gamma_i+\max_i\gamma_i,\\
 R_d&=2w_*,&
 R_l&=\ell(n)+\ell(1+mn^2)+2\gamma_*,\\
 A_d&=bt,& A_l&=b(\mu+\ell(b)),\\
 E_l&=A_l+\ell(b).&&
\end{align*}
\begin{align*}
 u_i&=d_i(d_i-1)w_i,&
 v_i&=d_i(d_i-1)\gamma_i+
             \ell\bigl((2d_i-1)!\,d_i^{d_i}\bigr),\\
 J_i&=b(A_d+u_i),&
 L_i&=b(E_l+v_i+2+\ell(b)),\\
 D_d&=2\sum_iJ_i,&D_l&=2\sum_iL_i,\\
 K_d&=SR_d+2mbA_d+D_d,\\
 K_l&=S(1+R_l)+mb(1+2A_l)+m+D_l,\\
 C_d&=R_d+K_d,\qquad C_l=3+R_l+K_l,\\
 W&=2A_d+D_d+mC_d+\max_i w_i,\\
 B_l&=A_l+D_l+E_l+\ell(m+1)+mC_l+\max_i\gamma_i.
\end{align*}
Subscripts $d,l$ mean degree and logarithmic coefficient-norm
bounds, respectively; $B_l$ is a number, not the algebra $B$.

\begin{theorem}\label{thm:bounds}
The explicit polynomial \eqref{eq:P} satisfies
\begin{align}
 \deg P&\le 2bn(W+1),\label{eq:degree}\\
 \log_2\|P\|_1&\le
 b\bigl(2n(B_l+\ell(n)+1)+\ell(b+1)\bigr).\label{eq:height}
\end{align}
Its coefficient height is at most $\|P\|_1$.
\end{theorem}
\begin{proof}
An entry in the multiplication matrix of an element of weight
$w$ has degree at most $w+$ column basis weight $-$ row basis
weight. The analogous logarithmic norm bound uses the $\gamma$
weights. Thus the entries of all $M_{\theta_i}$ have degree at
most $w_*$ and coefficient norm at most $2^{\gamma_*}$.
This proves the bounds $R_d,R_l$ for \eqref{eq:root-bound}.

In a principal minor, the row and column weights cancel.
If multiplication by an element has weight at most $u$ and
operator norm at most $2^v$ in an algebra of rank $h$, the
$j$th characteristic coefficient has degree at most $ju$ and
norm at most $\binom hj2^{jv}\le2^{j(v+\ell(h))}$.
Indeed one can expand the determinant in column sums after
the diagonal similarity defining the weighted norm.
In particular $a_j$ has degree at most $jt$ and logarithmic
norm at most $j(\mu+\ell(b))$.

The cancellation of the constant term in
$E=c-h_r(T)$ leaves exactly the $r$ terms with indices
$0,\ldots,r-1$. Its weight is at most $A_d$ and its
operator norm at most $2^{E_l}$. The native discriminant
is weighted homogeneous of weight $d_i(d_i-1)$ in its
roots. Its Sylvester determinant has at most $(2d_i-1)!$
terms, with at most $d_i$ derivative coefficients, each
contributing a factor at most $d_i$. This gives $u_i,v_i$.
The element $E\delta_i+c-E$ has weight at most $A_d+u_i$
and logarithmic operator norm at most $E_l+v_i+2$.
The determinant estimate just given bounds its norm by
$J_i,L_i$; the extra $\ell(b)$ in $L_i$ is harmless.

Products and sums of polynomial norms now give the displayed
bounds for $D,K,C$. For example
$\log_2\|1+c^2\|_1\le1+2A_l$ and
$\log_2\|1+D_i^2\|_1\le1+2L_i$.
Equation~\eqref{eq:beta} has weight at most $W$ and
operator norm at most $2^{B_l}$. Applying the same
principal-minor estimate to \eqref{eq:H} yields
\[
 \deg H_r\le n(W+1),\qquad
 \log_2\|H_r\|_1\le n(B_l+\ell(n)+1).
\]
These bounds dominate those for each $a_j$ in the applications,
and in the displayed general formulas since $m,n,b\ge1$.
Squaring and summing at most $b+1$ terms bounds $Q_r$.
Summing the bounds over the $b$ factors proves
\eqref{eq:degree} and \eqref{eq:height}.
\end{proof}

\section{The specified arithmetic constructions}

The guarded sources in \cite{Source27,Source29} have respectively
$26$ and $28$ retained base variables, including the input.
Each has nine independent scalar monic quartics for its point
charts, followed by respectively three and four monic quadratic
relations. Exact inspection of every coefficient dependency shows
that each quadratic radicand belongs to the point algebra and that
each pole generator occurs to exponent exactly one in the complete
guard. The other guard factors belong to the point algebra.
There are no residual equations or additional indexed guards.

The point-chart proof in \cite{Source27,Source29} shows separability
at every rational base point. More explicitly, its nonvanishing
follows from the nonsquare quadratic norm and the anisotropy over
$\Q$ of the transformed cubic. The discriminant is the stated
product
$2^{15}B_0^2(2U_0^2-V_0^2)^4(U_0^2-2V_0^2)^2$;
all factors are nonzero at every rational chart input.
Thus the hypotheses of Theorem~\ref{thm:main} hold on every
rational base point, without replacing a branchwise guard by
the stronger requirement that its norm be nonzero.

\begin{center}
\begin{tabular}{lrr}
\toprule
Quantity & $27$ variables & $29$ variables\\
\midrule
$m$ & $12$ & $13$\\
$b$ & $2^{18}$ & $2^{18}$\\
$n$ & $2^{21}$ & $2^{22}$\\
$w_G$ & $4\,423\,559$ & $6\,789\,625$\\
$\eta$ at $L=137731$ & $6\,998\,669\,574$ & $7\,005\,250\,004$\\
$S$ & $29\,884\,416$ & $30\,408\,704$\\
Number of factors & $262\,144$ & $262\,144$\\
\bottomrule
\end{tabular}
\end{center}
The point generator weights are all $6$ and their logarithmic
operator bounds are all $14$. The remaining weights are
$1568138,1935231,116318$ in the first source, and
$1568138,1935231,1935231,116318$ in the second.
Every integer recurrence and full logarithmic generator list is
recorded in the accompanying coefficient specifications.

\begin{corollary}\label{cor:27}
Under the arithmetic premises of the twenty-seven-variable source,
the construction gives a computable family in $27$ total variables
such that every sufficiently large odd fixed coefficient exponent
defines $\Z$ over $\Q$, and
\[
 \deg P\le401086041810161294689450143514624<2^{109}.
\]
The sufficient coefficient threshold is ineffective.
\end{corollary}
\begin{proof}
Theorem~\ref{thm:main} preserves the source projection exactly.
Its native degree weights are independent of the fixed coefficient
exponent. Substitute them into \eqref{eq:degree} and use the
source arithmetic theorem, with its unchanged threshold.
\end{proof}

\begin{corollary}\label{cor:29}
Under the arithmetic premises of the explicit twenty-nine-variable
source, the construction with fixed exponent $L=137731$ defines
$\Z$ over $\Q$ in $29$ total variables and satisfies
\[
 \deg P\le1440544063570503434967765431091200<2^{111},
\]
and
\[
 \log_2 H(P)\le
 1486299933480681227980829420111527936<2^{121}.
\]
\end{corollary}
\begin{proof}
Apply Theorem~\ref{thm:main} to the fixed effective source and
substitute its integer data into \eqref{eq:degree} and
\eqref{eq:height}. No arithmetic parameter is enlarged.
\end{proof}

The numerical $27$-variable member at the same exponent has the
certified bound $\log_2 H(P)<2^{119}$, but that exponent is only
an algebraic reproduction fixture for this family. It is not
certified to exceed its ineffective defining threshold.
None of the algebraic calculations replaces a source arithmetic
premise by a computational claim.

\section{Exact specification and verification}

The program \nolinkurl{build_active_specification.py} reads the
pinned integer coefficient circuits, verifies the dependency and
guard-exponent assertions, recomputes the weighted companion
bounds, and writes the complete formulas \eqref{eq:root-bound}--%
\eqref{eq:P} with their exact finite ranges.
The independent checker reconstructs the integer bounds and
verifies the hashes, the structural assertions, and the source
projection scope. Small exact algebra examples additionally cover
repeated nonzero $T$-values, inactive nilpotents, absence of rational
characters, and the empty active part. They illustrate specializations
of the theorem; the general proof is given above.

The rank-sized matrices and the expanded monomial list are not
materialized. An indexed expression using determinants, coefficient
extraction, integer sums, and finite products specifies every
coefficient without introducing an existential variable.

\begin{thebibliography}{9}
\bibitem{Source27}
\emph{A twenty-seven-variable polynomial definition of the integers},
the guarded arithmetic source and integer coefficient specification,
in the accompanying repository, \nolinkurl{papers/few_variables27}.
\bibitem{Source29}
\emph{An explicit twenty-nine-variable polynomial definition of the integers},
the effective guarded arithmetic source and integer coefficient
specification, in the accompanying repository,
\nolinkurl{papers/few_variables}.
\bibitem{Weighted}
\emph{Weighted bounds for eliminating finite algebraic choices},
the weighted companion-column and principal-minor estimates,
in the accompanying repository, \nolinkurl{papers/finite_algebra}.
\end{thebibliography}
\end{document}
